% MATLAB Script to Plot Water Height in Emptying Tank
 % Parameters
 rho = 1000; % density ( kg / m ^3)
 D = 1; % diameter ( m )
 h0 = 2; % initial water height ( m )
 g = 9.81; % gravitational acceleration ( m / s ^2)
 C = 0.001; % valve constant ( m ^3/ s / Pa )
 % Derived constants
 k = 4 * C * rho * g / ( pi * D ^2);

 % Time to reach half - empty height
 t_half = ( pi * D ^2 * log (2)) / (4 * C * rho * g );

 % Time to reach 0.001 h0 ( effectively empty )
 t_empty = ( pi * D ^2 * log (1000)) / (4 * C * rho * g );

 % Time vector from 0 to t_empty
 t = linspace (0 , t_empty , 1000);

 % Water height as a function of time
 h = h0 * exp ( - k * t );

 % Average height between t =0 and t_half
 h_avg = h0 / (2 * log (2)); % Average height

 % Plotting
 figure ;
 plot (t , h , ’b - ’ , ’ LineWidth ’ , 2); % Water height vs time
 hold on ;
 yline ( h_avg , ’r - - ’ , ’ LineWidth ’ , 2); % Average height line
 xlabel ( ’ Time ␣ ( s ) ’ );
 ylabel ( ’ Water ␣ Height ␣ ( m ) ’ );
 title ( ’ Water ␣ Height ␣ in ␣ an ␣ Emptying ␣ Tank ’ );
 legend ( ’ Water ␣ Height ’ , ’ Average ␣ Height ␣ between ␣ t =0 ␣ and ␣ t_ {1/2} ’ ,...
 ’ Location ’ , ’ Best ’ );
 grid on ;
 hold off ;


